\section{引言：连接图谱并不等于会学习的大脑}
\label{sec:intro}
成年果蝇脑的突触级连接组为计算建模提供了少见的解剖约束：我们不必从随机连接开始，而可以使用真实神经元之间的连接方向及接触数量\citep{dorkenwald2024}。Shiu 等人据此构建全脑泄漏积分发放模型，展示了味觉相关活动、感觉通路之间的相互作用以及触角梳理相关回路的可预测性\citep{shiu2024}。然而，单次刺激能够沿正确通路产生响应，不代表该模型已经拥有跨刺激记忆。

习惯化适合检验这个差别。与复杂任务训练相比，它不要求语言、显式奖赏或长期参数优化，却要求系统的当前响应依赖最近的刺激历史\citep{rankin2009}。如果一个模型每次收到同样输入都作出同样响应，它可以完成感觉传递，但没有在相应协议下表现出反应递减。相反，若引入某个慢变量后出现下降，也不能仅凭一条曲线确定动物使用了相同机制。

本研究提出两个递进问题：\textbf{（Q1）}原始固定连接的全脑动力学，在重复 JO-CE 刺激和足够长的间隔下，是否产生跨次响应递减与恢复？\textbf{（Q2）}如果没有，加入局部、可恢复的短时突触资源状态，能否在保留真实连接结构的条件下产生这些特征？本轮并未直接检验更强的\textbf{（Q3）}：真实连接组是否比简单动力学模型或随机结构具有不可替代的预测优势。

本报告的贡献是建立连续状态、输入严格配对且可审计的计算基线，记录一个明确的阴性结果和一个有局限的充分性结果。\textbf{短时突触抑制不是本研究新提出的生物学机制；“最小”指本轮只加入一种候选慢机制，不是已经证明全局数学最小性。}

\section{背景与概念：给没有神经科学基础的读者}
\label{sec:background}
\subsection{神经元、突触与连接组分别是什么}
神经元可以用膜电位表示其电活动状态。真实神经元的机制很复杂，本模型把它简化为一个会积累输入、同时向静息状态泄漏的单元。膜电位超过阈值时，单元发出一次离散事件，称为脉冲或动作电位；随后状态被重置，并短暂进入不应期。

突触是一个神经元影响另一个神经元的连接。\emph{兴奋性}影响在本模型中增加后续发放的倾向；\emph{抑制性}影响则减少这种倾向。它们不是“好”和“坏”的连接：抑制有助于选择响应、稳定回路，并可能参与习惯化。两个神经元之间可以有多个解剖突触，本模型把这些接触聚合为一条带权有向边。因此，\textbf{神经元数、聚合边数和解剖突触数是三个不同的数量}。

连接组回答“哪些细胞连接到哪些细胞”，但通常不直接提供所有膜时间常数、突触释放动力学或动物当时的内部状态。把连接组转化为一个可运行模型，还必须作出关于这些动力学的假设。解剖连接也不应直接等同于被测动物完整的生前记忆。

\begin{primer}{地图、交通规则与实时车流}
可以把连接组类比为道路地图，神经元和突触方程类比为交通规则，随时间改变的膜电位、脉冲和资源量则类似实时车流。有了地图，不代表交通规则已经确定；规则固定，也不代表车流没有历史依赖。这个类比只帮助理解结构与状态的区别，不是大脑与道路的机制等同。
\end{primer}

\subsection{为什么“不改长期权重”也可能保留记忆}
本研究区分三类量。第一类是\emph{结构}，如边是否存在、突触数量和连接方向；第二类是\emph{参数}，如阈值、时间常数和总体权重尺度；第三类是\emph{动态状态}，如当前膜电位、剩余突触驱动或可用资源。即使结构和参数保持不变，状态仍可随过去的刺激而变化，并影响下一次输出。

与典型机器学习中的优化不同，本轮没有梯度下降、训练损失或逐边权重拟合。M1 中的有效传递强度确实会随资源改变，因此它包含短时可塑性；但这不等于做了长期权重训练。本文所说的“12 次训练刺激”是习惯化实验的暴露序列，不是运行一个机器学习优化器。

原始 LIF 网络本来就有动态状态和循环连接。不能从“没有专门的可塑性规则”推出“固定权重网络绝不可能习惯化”。循环网络也可能形成慢活动模式。我们报告的是实际测试范围内的阴性结果，而不是对所有固定权重网络的否定。

\subsection{习惯化、感觉适应与疲劳为什么不能混为一谈}
习惯化通常指重复刺激导致的反应减弱，且不能简单归因于感觉器官越来越不灵敏或运动器官失去执行能力\citep{rankin2009}。观察到输出下降至少有四种可能：输入真的变弱了；外围感觉器官适应了；中枢处理保留了历史；输出通路失能了。因此，“越来越少发放”只是现象，不是机制诊断。

经典研究列出了反应递减、休息恢复、频率依赖、刺激特异性和去习惯化等特征。所有系统不必同时满足全部特征，但研究者应清楚说明检验了哪些。尤其要区分两种测试：训练 A 后对 B 仍有响应，属于特异性或泛化问题；插入 B 后对原刺激 A 的响应恢复，才是去习惯化测试。后者还需要排除休息本身带来的恢复以及普遍敏化。

本轮把相同感觉脉冲直接回放给模型，控制了刺激编码的变化；同时检测了神经读出的休息恢复。它没有真实感觉器官和完整身体，因此\textbf{不能声称已在生物实验意义上排除全部感觉适应或运动疲劳解释}。

\subsection{相关工作如何约束模型选择}
最小动力学研究表明，一个会积累并衰减的内部状态，加上非线性读出，就可能生成多项习惯化特征；因此，全脑规模并不是产生下降和恢复曲线的必要前提\citep{smart2024}。这种认识提高了连接组研究的验收标准：真正有价值的下一步是预测刺激变化和神经干预，而不只是解释一条已经看过的曲线。

生物实验还提供了不同的候选机制。果蝇嗅觉习惯化研究支持局部 GABA 能抑制性传递的可塑性；在相应实验中，短期习惯化涉及约 30 分钟的气味暴露及分钟级恢复\citep{das2011}。味觉研究支持抑制性神经元相关机制，并进一步研究了新刺激如何解除已经形成的习惯化\citep{paranjpe2012,trisal2022}。这些结果支持把慢抑制和去抑制列为竞争模型，但不能直接证明 JO-CE 也依赖相同机制。

视觉 light-off 惊跳和其他果蝇范式具有丰富的习惯化研究基础\citep{engel2009}；近期个体动力学建模进一步说明，低维状态模型具有相当的解释能力\citep{boon2026}。本轮没有把任意视觉激活当作 light-off，也没有把秒级神经输出替代分钟级嗅觉或味觉行为。

\subsection{为什么选 JO-CE，而不是继续局限于糖刺激}
Johnston 器官是果蝇触角中的机械感觉器官。JO-CE 是其中一组感觉神经元；aBN1 是触角梳理相关中间神经元，aDN1 和 aDN2 是相关下行神经元。原全脑研究已经给出这些细胞的可用标识，且 JO-CE 激活 aBN1 的预测得到生理验证\citep{shiu2024}。这为正常反应提供了依据，允许先讨论“重复之后发生什么”。

选择该通路是本轮可操作性与科学可证伪性的折中，并非认定它比嗅觉或视觉更重要。输入被定义为神经脉冲，不等价于已校准的声音、振动波形或触角触碰。JO-F 在原模型中不能可靠激活 aBN1，因此不能简单把 JO-F 当作一个可直接比较的“新刺激 B”。最重要的是，aBN1 的活动只是神经读出，不是果蝇梳理动作本身。
